\documentclass[a4paper,12pt,twoside]{memoir}

% Castellano
\usepackage[spanish,es-tabla]{babel}
\selectlanguage{spanish}
\usepackage[utf8]{inputenc}
\usepackage[T1]{fontenc}
\usepackage{lmodern} % Scalable font
\usepackage{microtype}
\usepackage{placeins}

\RequirePackage{booktabs}
\RequirePackage[table]{xcolor}
\RequirePackage{xtab}
\RequirePackage{multirow}

% Links
\PassOptionsToPackage{hyphens}{url}\usepackage[colorlinks]{hyperref}
\hypersetup{
	allcolors = {red}
}

% Ecuaciones
\usepackage{amsmath}

% Rutas de fichero / paquete
\newcommand{\ruta}[1]{{\sffamily #1}}

% Párrafos
\nonzeroparskip

% Huérfanas y viudas
\widowpenalty100000
\clubpenalty100000

% Imágenes

% Comando para insertar una imagen en un lugar concreto.
% Los parámetros son:
% 1 --> Ruta absoluta/relativa de la figura
% 2 --> Texto a pie de figura
% 3 --> Tamaño en tanto por uno relativo al ancho de página
\usepackage{graphicx}
\newcommand{\imagen}[3]{
	\begin{figure}[!h]
		\centering
		\includegraphics[width=#3\textwidth]{#1}
		\caption{#2}\label{fig:#1}
	\end{figure}
	\FloatBarrier
}

% Comando para insertar una imagen sin posición.
% Los parámetros son:
% 1 --> Ruta absoluta/relativa de la figura
% 2 --> Texto a pie de figura
% 3 --> Tamaño en tanto por uno relativo al ancho de página
\newcommand{\imagenflotante}[3]{
	\begin{figure}
		\centering
		\includegraphics[width=#3\textwidth]{#1}
		\caption{#2}\label{fig:#1}
	\end{figure}
}

% El comando \figura nos permite insertar figuras comodamente, y utilizando
% siempre el mismo formato. Los parametros son:
% 1 --> Porcentaje del ancho de página que ocupará la figura (de 0 a 1)
% 2 --> Fichero de la imagen
% 3 --> Texto a pie de imagen
% 4 --> Etiqueta (label) para referencias
% 5 --> Opciones que queramos pasarle al \includegraphics
% 6 --> Opciones de posicionamiento a pasarle a \begin{figure}
\newcommand{\figuraConPosicion}[6]{%
  \setlength{\anchoFloat}{#1\textwidth}%
  \addtolength{\anchoFloat}{-4\fboxsep}%
  \setlength{\anchoFigura}{\anchoFloat}%
  \begin{figure}[#6]
    \begin{center}%
      \Ovalbox{%
        \begin{minipage}{\anchoFloat}%
          \begin{center}%
            \includegraphics[width=\anchoFigura,#5]{#2}%
            \caption{#3}%
            \label{#4}%
          \end{center}%
        \end{minipage}
      }%
    \end{center}%
  \end{figure}%
}

%
% Comando para incluir imágenes en formato apaisado (sin marco).
\newcommand{\figuraApaisadaSinMarco}[5]{%
  \begin{figure}%
    \begin{center}%
    \includegraphics[angle=90,height=#1\textheight,#5]{#2}%
    \caption{#3}%
    \label{#4}%
    \end{center}%
  \end{figure}%
}
% Para las tablas
\newcommand{\otoprule}{\midrule [\heavyrulewidth]}
%
% Nuevo comando para tablas pequeñas (menos de una página).
\newcommand{\tablaSmall}[5]{%
 \begin{table}
  \begin{center}
   \rowcolors {2}{gray!35}{}
   \begin{tabular}{#2}
    \toprule
    #4
    \otoprule
    #5
    \bottomrule
   \end{tabular}
   \caption{#1}
   \label{tabla:#3}
  \end{center}
 \end{table}
}

%
% Nuevo comando para tablas pequeñas (menos de una página).
\newcommand{\tablaSmallSinColores}[5]{%
 \begin{table}[H]
  \begin{center}
   \begin{tabular}{#2}
    \toprule
    #4
    \otoprule
    #5
    \bottomrule
   \end{tabular}
   \caption{#1}
   \label{tabla:#3}
  \end{center}
 \end{table}
}

\newcommand{\tablaApaisadaSmall}[5]{%
\begin{landscape}
  \begin{table}
   \begin{center}
    \rowcolors {2}{gray!35}{}
    \begin{tabular}{#2}
     \toprule
     #4
     \otoprule
     #5
     \bottomrule
    \end{tabular}
    \caption{#1}
    \label{tabla:#3}
   \end{center}
  \end{table}
\end{landscape}
}

%
% Nuevo comando para tablas grandes con cabecera y filas alternas coloreadas en gris.
\newcommand{\tabla}[6]{%
  \begin{center}
    \tablefirsthead{
      \toprule
      #5
      \otoprule
    }
    \tablehead{
      \multicolumn{#3}{l}{\small\sl continúa desde la página anterior}\\
      \toprule
      #5
      \otoprule
    }
    \tabletail{
      \hline
      \multicolumn{#3}{r}{\small\sl continúa en la página siguiente}\\
    }
    \tablelasttail{
      \hline
    }
    \bottomcaption{#1}
    \rowcolors {2}{gray!35}{}
    \begin{xtabular}{#2}
      #6
      \bottomrule
    \end{xtabular}
    \label{tabla:#4}
  \end{center}
}

%
% Nuevo comando para tablas grandes con cabecera.
\newcommand{\tablaSinColores}[6]{%
  \begin{center}
    \tablefirsthead{
      \toprule
      #5
      \otoprule
    }
    \tablehead{
      \multicolumn{#3}{l}{\small\sl continúa desde la página anterior}\\
      \toprule
      #5
      \otoprule
    }
    \tabletail{
      \hline
      \multicolumn{#3}{r}{\small\sl continúa en la página siguiente}\\
    }
    \tablelasttail{
      \hline
    }
    \bottomcaption{#1}
    \begin{xtabular}{#2}
      #6
      \bottomrule
    \end{xtabular}
    \label{tabla:#4}
  \end{center}
}

%
% Nuevo comando para tablas grandes sin cabecera.
\newcommand{\tablaSinCabecera}[5]{%
  \begin{center}
    \tablefirsthead{
      \toprule
    }
    \tablehead{
      \multicolumn{#3}{l}{\small\sl continúa desde la página anterior}\\
      \hline
    }
    \tabletail{
      \hline
      \multicolumn{#3}{r}{\small\sl continúa en la página siguiente}\\
    }
    \tablelasttail{
      \hline
    }
    \bottomcaption{#1}
  \begin{xtabular}{#2}
    #5
   \bottomrule
  \end{xtabular}
  \label{tabla:#4}
  \end{center}
}



\definecolor{cgoLight}{HTML}{EEEEEE}
\definecolor{cgoExtralight}{HTML}{FFFFFF}

%
% Nuevo comando para tablas grandes sin cabecera.
\newcommand{\tablaSinCabeceraConBandas}[5]{%
  \begin{center}
    \tablefirsthead{
      \toprule
    }
    \tablehead{
      \multicolumn{#3}{l}{\small\sl continúa desde la página anterior}\\
      \hline
    }
    \tabletail{
      \hline
      \multicolumn{#3}{r}{\small\sl continúa en la página siguiente}\\
    }
    \tablelasttail{
      \hline
    }
    \bottomcaption{#1}
    \rowcolors[]{1}{cgoExtralight}{cgoLight}

  \begin{xtabular}{#2}
    #5
   \bottomrule
  \end{xtabular}
  \label{tabla:#4}
  \end{center}
}



\graphicspath{ {./img/} }

% Capítulos
\chapterstyle{bianchi}
\newcommand{\capitulo}[2]{
	\setcounter{chapter}{#1}
	\setcounter{section}{0}
	\setcounter{figure}{0}
	\setcounter{table}{0}
	\chapter*{\thechapter.\enskip #2}
	\addcontentsline{toc}{chapter}{\thechapter.\enskip #2}
	\markboth{#2}{#2}
}

% Apéndices
\renewcommand{\appendixname}{Apéndice}
\renewcommand*\cftappendixname{\appendixname}

\newcommand{\apendice}[1]{
	%\renewcommand{\thechapter}{A}
	\chapter{#1}
}

\renewcommand*\cftappendixname{\appendixname\ }

% Formato de portada
\makeatletter
\usepackage{xcolor}
\newcommand{\tutor}[1]{\def\@tutor{#1}}
\newcommand{\course}[1]{\def\@course{#1}}
\definecolor{cpardoBox}{HTML}{E6E6FF}
\def\maketitle{
  \null
  \thispagestyle{empty}
  % Cabecera ----------------
\noindent\includegraphics[width=\textwidth]{cabecera}\vspace{1cm}%
  \vfill
  % Título proyecto y escudo informática ----------------
  \colorbox{cpardoBox}{%
    \begin{minipage}{.8\textwidth}
      \vspace{.5cm}\Large
      \begin{center}
      \textbf{TFG del Grado en Ingeniería Informática}\vspace{.6cm}\\
      \textbf{\LARGE\@title{}}
      \end{center}
      \vspace{.2cm}
    \end{minipage}

  }%
  \hfill\begin{minipage}{.20\textwidth}
    \includegraphics[width=\textwidth]{escudoInfor}
  \end{minipage}
  \vfill
  % Datos de alumno, curso y tutores ------------------
  \begin{center}%
  {%
    \noindent\LARGE
    Presentado por \@author{}\\ 
    en Universidad de Burgos --- \@date{}\\
    Tutores: \@tutor{}\\
  }%
  \end{center}%
  \null
  \cleardoublepage
  }
\makeatother

\newcommand{\nombre}{Rodrigo Grande Muñoz} %%% cambio de comando

% Datos de portada
\title{VRRecoveryGym: entorno de realidad virtual para personas con problemas motrices}
\author{\nombre}
\tutor{Álvar Arnaiz González, Héctor Royo Concepción y Luis Fuente Ibáñez}
\date{\today}

\begin{document}

\maketitle


\newpage\null\thispagestyle{empty}\newpage


%%%%%%%%%%%%%%%%%%%%%%%%%%%%%%%%%%%%%%%%%%%%%%%%%%%%%%%%%%%%%%%%%%%%%%%%%%%%%%%%%%%%%%%%
\thispagestyle{empty}


\noindent\includegraphics[width=\textwidth]{cabecera}\vspace{1cm}

\noindent D. Álvar Arnaiz González, profesor del departamento de Ingeniería Informática, área de Lenguajes y Sistemas Informáticos.

\noindent Expone:

\noindent Que el alumno D. \nombre, con DNI 72901698R, ha realizado el Trabajo final de Grado en Ingeniería Informática titulado <<VRRecoveryGym: entorno de realidad virtual para personas con problemas motrices>>. 

\noindent Y que dicho trabajo ha sido realizado por el alumno bajo la dirección del que suscribe, en virtud de lo cual se autoriza su presentación y defensa.

\begin{center} %\large
En Burgos, {\large \today}
\end{center}

\vfill\vfill\vfill

% Author and supervisor
\begin{minipage}{0.45\textwidth}
\begin{flushleft} %\large
Vº. Bº. del Tutor:\\[2cm]
D. Álvar Arnaiz González
\end{flushleft}
\end{minipage}
\hfill
\begin{minipage}{0.45\textwidth}
\begin{flushleft} %\large
Vº. Bº. del co-tutor:\\[2cm]
D. Héctor Royo Concepción
\end{flushleft}
\end{minipage}
\vfill
\begin{minipage}{0.45\textwidth}
\begin{flushleft} %\large
Vº. Bº. del co-tutor:\\[2cm]
D. Luis Fuente Ibáñez
\end{flushleft}
\end{minipage}
\hfill

\vfill

% para casos con solo un tutor comentar lo anterior
% y descomentar lo siguiente
%Vº. Bº. del Tutor:\\[2cm]
%D. nombre tutor


\newpage\null\thispagestyle{empty}\newpage




\frontmatter

% Abstract en castellano
\renewcommand*\abstractname{Resumen}
\begin{abstract}
El objetivo de este proyecto es desarrollar un \textit {entorno virtual} en el que los usuarios puedan realizar estos ejercicios desde casa mediante herramientas de realidad virtual. Esta solución permitirá no solo la realización remota de los ejercicios, sino también un seguimiento preciso de la actividad de cada usuario, registrando parámetros que puedan ser consultados en remoto por profesionales del sector.


Este proyecto se realiza en colaboración con la fundación ASPAYM y HP, cuya misión es mejorar la calidad de vida de las personas con lesión medular y discapacidad física grave. Con el objetivo de favorecer su recuperación, la fundación ha desarrollado una serie de puzles impresos en 3D que fomentan movimientos específicos, como la extensión del codo, la abducción del brazo y el giro de la muñeca. Estos ejercicios han demostrado ser efectivos, pero requieren la presencia física de los usuarios en la asociación y la supervisión de un profesional.

Para el desarrollo de la aplicación se utilizará el motor Unreal Engine, empleando tanto C++ como \textit{Blueprint} para la programación de la lógica del sistema.

Se busca ofrecer una solución innovadora y accesible que facilite la rehabilitación motriz, mejorando la autonomía de los usuarios y optimizando el seguimiento terapéutico a través de la tecnología de realidad virtual.
\end{abstract}

\renewcommand*\abstractname{Descriptores}
\begin{abstract}
puzzles, \textit {entorno virtual}, seguimiento remoto, rehabilitación motriz, Unreal Engine
\end{abstract}

\clearpage

% Abstract en inglés
\renewcommand*\abstractname{Abstract}
\begin{abstract}
The goal of this project is to develop a  \textit {virtual environment} where users can perform these exercises from home using virtual reality tools. This solution will not only enable remote execution of exercises but also allow for precise tracking of each user's activity, recording parameters that can be remotely accessed by healthcare professionals.


This project is carried out in collaboration with the ASPAYM Foundation and HP, whose mission is to improve the quality of life of people with spinal cord injuries and severe physical disabilities. In order to support their recovery, the foundation has developed a series of 3D-printed puzzles designed to encourage specific movements, such as elbow extension, arm abduction, and wrist rotation. These exercises have proven to be effective, but they require users to be physically present at the association and supervised by a professional.

The application will be developed using the Unreal Engine, utilizing both C++ and \textit{Blueprint} for system logic programming.

This project aims to provide an innovative and accessible solution that facilitates motor rehabilitation, enhances users' autonomy, and optimizes therapeutic monitoring through virtual reality technology.
\end{abstract}

\renewcommand*\abstractname{Keywords}
\begin{abstract}
puzzles,  \textit {virtual environment}, remote monitoring, motor rehabilitation, Unreal Engine
\end{abstract}

\clearpage

% Indices
\tableofcontents

\clearpage

\listoffigures

\clearpage

\listoftables
\clearpage

\mainmatter
\capitulo{1}{Introducción}


Vivimos en una época donde la tecnología se ha convertido en una herramienta clave para mejorar la calidad de vida de las personas. En particular, el desarrollo de soluciones tecnológicas aplicadas al ámbito de la salud ha abierto nuevas posibilidades para abordar desafíos de rehabilitación física mediante enfoques innovadores. En este contexto, la realidad virtual (VR) ha demostrado un gran potencial para motivar y asistir a pacientes en procesos de recuperación motriz, especialmente en colectivos con discapacidades físicas que requieren terapias continuadas y personalizadas.

Actualmente, se utilizan puzles físicos impresos en 3D para promover movimientos específicos que permiten la rehabilitación de estas personas. Sin embargo, este enfoque presencial limita la autonomía de los usuarios y requiere supervisión constante por parte de profesionales.

El presente proyecto busca trasladar estos ejercicios a un entorno virtual accesible desde el domicilio del usuario. Mediante el uso de cascos de realidad virtual y controladores (o \textit{hand tracking} en un futuro), se desarrollarán actividades interactivas que simulan los puzles físicos. Además, se recopilarán métricas clave sobre la actividad motriz y el progreso del usuario, que podrán ser consultadas por profesionales de forma remota mediante un \textit{CSV} exportado con todos los datos relevantes sobre el progreso del paciente en los puzles.

El desarrollo se llevará a cabo utilizando Unreal Engine, combinando programación en C++ y \textit{Blueprint}, así como recursos 3D entregados, ya realizados por la fundación, y el \textit{marketplace} de Unreal. El proyecto se beneficia del asesoramiento de los profesionales de ASPAYM para garantizar la validez terapéutica del sistema.

\capitulo{2}{Objetivos del proyecto}

Este apartado explica de forma precisa y concisa cuales son los objetivos que se persiguen con la realización del proyecto. Se puede distinguir entre los objetivos marcados por los requisitos del software a construir y los objetivos de carácter técnico que plantea a la hora de llevar a la práctica el proyecto.

\capitulo{3}{Conceptos teóricos}
En esta sección se presentan los fundamentos teóricos necesarios para comprender el desarrollo y la implementación del presente Trabajo de Fin de Grado. Se abordan las principales definiciones, teorías y modelos que sustentan la investigación, así como las tecnologías y herramientas empleadas. 
\section{Realidad virtual}
Se entiende por realidad virtual (RV) la tecnología que permite crear entornos digitales tridimensionales, interactivos y en tiempo real, de forma que el usuario perciba la sensación de “estar presente” en dicho entorno. Este grado de inmersión y “telepresencia” se consigue mediante la combinación de hardware (cascos HMD, sistemas de tracking, interfaces hápticas) y software (motores gráficos, algoritmos de simulación) que actualizan la escena virtual en función de los movimientos y acciones del usuario ~~\cite{steuer92,sherman2002}.

Dentro del estudio de la RV suelen distinguirse tres características fundamentales:

\textbf{Inmersión:} grado en que el sistema aísla y sustituye estímulos sensoriales reales por estímulos generados artificialmente, creando la impresión directa de “estar dentro” del entorno virtual ~~\cite{milgram94}.

\textbf{Presencia:} sensación subjetiva de “estar realmente” en el escenario virtual, más allá de la mera percepción sensorial ~~\cite{sherman2002,slater94}.

\textbf{Interactividad:} capacidad del sistema para registrar las acciones del usuario y modificar el entorno virtual de forma coherente y en tiempo real ~~\cite{steuer92,milgram94}.

Estas características hacen de la RV una herramienta poderosa en ámbitos tan diversos como la simulación de entrenamiento, la rehabilitación médica, la visualización arquitectónica o la investigación científica.

\subsection{Componentes de un sistema de RV}

Un sistema de RV típico se compone de los siguientes elementos~\cite{sherman2002,burdea2003}:
\begin{itemize}
  \item \textbf{Dispositivo de visualización (Head-Mounted Display, HMD):} muestra las imágenes estereoscópicas e incorpora sensores de orientación (giroscopios, acelerómetros).
  \item \textbf{Sistema de seguimiento (tracking):} rastrea la posición y orientación de la cabeza, manos u otros objetos mediante cámaras, emisores infrarrojos o sistemas magnéticos.
  \item \textbf{Dispositivos de interacción:} mandos, guantes hápticos o superficies táctiles que permiten manipular objetos virtuales.
  \item \textbf{Motor gráfico y software de simulación:} genera las escenas 3D, gestiona física, colisiones y comportamientos de los elementos del entorno.
\end{itemize}

\subsection{Taxonomía: Continuo de la virtualidad}

Milgram y Kishino propusieron el concepto de continuo de la virtualidad, que va desde la realidad completamente física hasta la realidad completamente virtual, englobando la realidad aumentada (AR) y la realidad mixta asistida ~~\cite{milgram94}:
\begin{itemize}
  \item \textbf{Realidad Aumentada (AR):} superposición de información virtual sobre el entorno real.
  \item \textbf{Realidad Mixta (MR):} fusión de elementos reales y virtuales que interactúan entre sí.
  \item \textbf{Realidad Virtual (RV):} reemplazo total del entorno real por uno virtual.
\end{itemize}

\subsection{Modalidades de uso y aplicaciones}

Según el grado de inmersión y el tipo de interacción, los sistemas de RV pueden clasificarse en:
\begin{itemize}
  \item \textbf{RV de escritorio:} utiliza pantallas convencionales y controladores estándar; menor inmersión, pero accesible.
  \item \textbf{RV inmersiva:} emplea HMDs y sistemas de tracking avanzado; alta inmersión y presencia ~~\cite{burdea2003}.
  \item \textbf{CAVEs:} espacios cerrados con pantallas proyectadas en paredes y suelo; múltiples usuarios simultáneos.
\end{itemize}

Las aplicaciones más difundidas incluyen simuladores de vuelo y conducción, terapias de exposición en salud mental, entrenamiento quirúrgico y entornos de aprendizaje inmersivo.

\section{Diseño de videojuegos para rehabilitación}

El uso de videojuegos para la rehabilitación física se enmarca dentro del concepto de \textit{serious games}, definido como aquellos juegos cuyo principal objetivo va más allá del entretenimiento, buscando también la educación, la formación o la mejora de la salud ~~\cite{david2016serious}.

\subsection{Principios de diseño}

El desarrollo de videojuegos para rehabilitación debe considerar aspectos tanto lúdicos como terapéuticos ~~\cite{greene2016rehab}. Entre los principales principios de diseño se encuentran:

\begin{itemize}
    \item \textbf{Interactividad:} permite al usuario manipular y explorar el entorno, facilitando el aprendizaje y la mejora de habilidades motoras.
    \item \textbf{\textit{Feedback} inmediato:} el sistema debe ofrecer respuestas instantáneas a las acciones del usuario, reforzando el progreso ~~\cite{cameirao2010impact}.
    \item \textbf{Dificultad progresiva:} las tareas deben adaptarse al nivel de capacidad del paciente, aumentando la complejidad de forma gradual para evitar la frustración ~~\cite{zimmerli2012virtual}.
    \item \textbf{Repetición y refuerzo positivo:} elementos clave en la rehabilitación física para promover la neuroplasticidad.~~\cite{holden2005virtual}.
\end{itemize}

\subsection{Gamificación y motivación}

La gamificación introduce mecánicas de juego (puntos, logros, niveles) que aumentan la motivación y el compromiso del usuario ~~\cite{sardi2017gamification}. En el caso del presente proyecto, la superación de puzzles y la obtención de logros sirven como incentivo para continuar con las sesiones terapéuticas.

\subsection{Accesibilidad y consideraciones éticas}

Los videojuegos de rehabilitación deben garantizar la accesibilidad a personas con diferentes grados de discapacidad. Además, se debe asegurar que el uso de la tecnología no sustituya la supervisión profesional, sino que actúe como complemento ~~\cite{rodriguez2020reha}.

\section{Teoría del aprendizaje motor}

La teoría del aprendizaje motor se basa en el modelo trifásico de Fitts y Posner \cite{fitts1967}, que describe la adquisición de habilidades motoras en tres etapas:

\begin{enumerate}
  \item \textbf{Fase cognitiva:} el aprendiz construye una representación mental de la tarea, explorando secuencias de movimiento y cometiendo errores frecuentes. La ejecución es lenta y requiere alta atención consciente \cite{schmidt2019}.
  \item \textbf{Fase asociativa:} se establecen mapeos más precisos entre estímulos y respuestas motoras. La variabilidad de la ejecución disminuye y se corrigen errores de forma más autónoma, consolidándose patrones de movimiento \cite{schmidt2019}.
  \item \textbf{Fase autónoma:} la habilidad está completamente automatizada; la ejecución es rápida, fluida y apenas requiere atención consciente, liberando recursos cognitivos para otras tareas \cite{fitts1967}.
\end{enumerate}

El \emph{feedback} inmediato es crítico para la consolidación de estos patrones. Winstein y Schmidt demostraron que una frecuencia reducida de conocimiento de resultados (\emph{knowledge of results}) optimiza el aprendizaje a largo plazo, evitando la dependencia excesiva de la retroalimentación externa \cite{winstein1991}. En entornos de realidad virtual, el \emph{feedback} inmediato y multimodal (visual, auditivo y háptico) permite detectar y corregir errores en tiempo real, reforzando los esquemas motores y favoreciendo la neuroplasticidad \cite{schmidt2019}. Además, la inmersión en RV incrementa la motivación y la adherencia al entrenamiento, potenciando la repetición deliberada y la automatización de la habilidad.

\section{Gamificación y motivación intrínseca}

La \emph{gamificación} consiste en la aplicación de elementos de diseño de juegos (p. ej., puntos, niveles, insignias, narrativa, retos) en contextos no lúdicos con el objetivo de aumentar la implicación, la motivación y la adherencia de los usuarios a una tarea \cite{deterding2011defining}. En el ámbito de la rehabilitación en RV, la gamificación permite transformar ejercicios repetitivos en experiencias más atractivas, favoreciendo la práctica continua y la mejora progresiva de las capacidades motoras y cognitivas.

La \emph{motivación intrínseca} se sustenta en la Teoría de la Autodeterminación de Deci y Ryan \cite{deci1985intrinsic}, que postula tres necesidades psicológicas básicas:
\begin{itemize}
  \item \textbf{Autonomía:} sensación de control y elección sobre la propia conducta.
  \item \textbf{Competencia:} percepción de eficacia y dominio en la realización de tareas.
  \item \textbf{Relación:} conexión social y sentido de pertenencia con otros participantes o terapeutas.
\end{itemize}
Cuando un sistema gamificado satisface estas necesidades, la motivación intrínseca aumenta, conduciendo a una práctica más sostenida y a mejores resultados terapéuticos.

Kapp \cite{kapp2012gamification} propone que la gamificación efectiva integra mecánicas (reglas del juego), dinámicas (comportamientos de los usuarios) y elementos estéticos (presentación), de modo que:
\begin{enumerate}
  \item \textbf{Mecánicas:} puntos, niveles, temporizadores y desafíos escalonados para construir progresión.
  \item \textbf{Dinámicas:} recompensas variables, retroalimentación inmediata y metas claras que fomentan la repetición deliberada.
  \item \textbf{Estética:} narrativa inmersiva y diseño coherente con la experiencia de usuario en RV.
\end{enumerate}
En entorno de RV, el \emph{feedback} inmediato (visual, auditivo y háptico) refuerza la sensación de competencia y contribuye a la consolidación de los aprendizajes motores y cognitivos.

\capitulo{4}{Técnicas y herramientas}

En esta sección se presentan algunas de las herramientas fundamentales empleadas en el desarrollo del proyecto.
\section{Herramientas}
\subsection{Unreal Engine}

Unreal Engine es un motor gráfico y de desarrollo de juegos en tiempo real desarrollado por Epic Games. Su arquitectura modular y su potente editor visual permiten diseñar, iterar y desplegar entornos 3D de alta fidelidad para aplicaciones que van desde videojuegos hasta simulaciones VR y AR. Para este proyecto se ha utilizado Unreal Engine 5, aprovechando sus siguientes características principales:

\begin{itemize}
  \item \textbf{Editor de niveles y blueprints:} interfaz WYSIWYG para construir escenas y lógica de juego mediante nodos, sin necesidad de programar directamente en C++.
  \item \textbf{Renderizado de última generación:} soporte para Nanite (geometría virtualizada) y Lumen (iluminación global dinámica), que permiten obtener gráficos realistas en tiempo real.
  \item \textbf{Soporte nativo de Realidad Virtual:} integración con SDKs de dispositivos HMD (Oculus, SteamVR, etc.), sistemas de tracking y controladores hápticos.
\end{itemize}

Para la integración de la interfaz de usuario en RV se ha empleado el sistema UMG (Unreal Motion Graphics), lo que ha permitido diseñar menús y HUDs directamente dentro del entorno 3D.

\subsection{Blueprints} 
Las Blueprints constituyen el sistema de scripting visual de Unreal Engine que permite definir la lógica de juego y comportamiento de objetos sin escribir código en C++ directamente. Basado en un paradigma de programación orientada a objetos y data flow, Blueprints ofrece una interfaz de nodos conectables que representan:

\textbf{Clases y Actores:} cada Blueprint es esencialmente una subclase de \textit{UObject} o \textit{AActor}, con variables (propiedades), funciones y eventos propios.

\textbf{Graph Editor:} el Event Graph agrupa nodos de eventos (p. ej., BeginPlay, Tick, entradas de usuario) con nodos de ejecución (funciones, macros, secuencias lógicas) que conectan flujos de datos y control.

\textbf{Construction Script:} un grafo especial que se ejecuta en el editor para inicializar propiedades y componentes antes del runtime.

\textbf{Typed Nodes:} Blueprints maneja múltiples tipos de nodos (“Event”, “Function”, “Macro”, “Interface”, “Custom Event”), cada uno con firma de entrada/salida, lo que facilita la reutilización de lógica.

\textbf{Compilación y VM:} al guardar un Blueprint, Unreal compila el grafo a un bytecode que corre en la Blueprint Virtual Machine, garantizando rendimiento y depuración (breakpoints, watch variables) similares a código nativo \cite{epicBlueprints}.

\subsection{Blender}

Blender es una suite de creación 3D de código abierto que incluye herramientas para modelado, escultura, texturizado y exportación de geometrías. En este proyecto, su uso se ha centrado exclusivamente en la preparación de los modelos 3D proporcionados en formato STL y su posterior exportación para Unreal Engine:

\begin{itemize}
  \item \textbf{Importación de STL:} carga de los ficheros STL originales, comprobación de integridad del modelo.
  \item \textbf{Ajuste de escala y orientación:} normalización de las dimensiones del modelo y alineación de los ejes para que coincidan con las unidades y el sistema de coordenadas de Unreal Engine.
  \item \textbf{Conversión y exportación:} exportación del modelo en formatos compatibles (FBX u OBJ), asegurando que las normales de las caras permanezcan orientadas correctamente y que no se pierda información geométrica al importar en Unreal Engine.
  
\end{itemize}
\capitulo{5}{Aspectos relevantes del desarrollo del proyecto}

Uno de los aspectos más significativos de este proyecto ha sido la oportunidad de trabajar en un entorno real con un cliente externo. En este caso, el cliente fue el responsable de la asociación ASPAYM, quien participó activamente en las reuniones de seguimiento del proyecto. Su implicación permitió establecer una comunicación directa y constante, aportando ideas, validando propuestas y sugiriendo modificaciones en función de las necesidades reales del colectivo al que va dirigido el sistema. Esta experiencia ha sido especialmente enriquecedora, ya que ha permitido integrar la perspectiva del usuario final en todas las fases del desarrollo, adaptando el enfoque inicial a una solución verdaderamente útil y accesible.

Otro punto destacable ha sido el proceso de aprendizaje en torno al desarrollo de aplicaciones en \textit{realidad virtual}, un campo completamente nuevo al inicio del trabajo. El uso de \textit{Unreal Engine}, junto con el desarrollo para dispositivos \textit{Meta Quest}, ha supuesto un reto técnico importante que ha requerido una formación autodidacta intensiva. A lo largo del proyecto, se ha adquirido experiencia en el diseño de entornos virtuales inmersivos, implementación de mecánicas interactivas y gestión de la lógica del juego en un entorno 3D, aplicando tanto programación en \textit{Blueprints} como nociones básicas de \textit{C++}.

Además, el uso de recursos 3D ha sido fundamental para el diseño de los puzles interactivos del sistema. Se han integrado modelos tridimensionales que no solo cumplen una función estética si no que ha requerido un trabajo minucioso de adaptación, configuración de colisiones y ajuste de físicas para garantizar una experiencia fluida, funcional y segura para el usuario.

Por ultimo, resulta crucial haber identificado desde el inicio las particularidades de los usuarios finales, muchos de los cuales presentan limitaciones motrices derivadas de su condición. Diseñar la jugabilidad pensando en ellas implicó adoptar mecánicas que reduzcan la necesidad de movimientos precisos y de alta velocidad, así como ofrecer múltiples modalidades de interacción—ya sea mediante gestos amplios y simplificados, menús accesibles o ayudas visuales y sonoras que guíen al jugador. Este enfoque centrado en la accesibilidad no solo mejora la experiencia de quienes tienen dificultades físicas, sino que también enriquece el sistema general al fomentar una interfaz intuitiva y adaptable, capaz de ajustarse a distintos niveles de destreza y garantizar una participación inclusiva y satisfactoria.
\capitulo{6}{Trabajos relacionados}

Este apartado sería parecido a un estado del arte de una tesis o tesina. En un trabajo final grado no parece obligada su presencia, aunque se puede dejar a juicio del tutor el incluir un pequeño resumen comentado de los trabajos y proyectos ya realizados en el campo del proyecto en curso. 

\capitulo{7}{Conclusiones y Líneas de trabajo futuras}

En este último apartado de la memoria se pretende exponer todo lo aprendido durante la realización del proyecto y mostrar ideas sobre mejoras para el presente proyecto.

\section{Conclusiones}

Debido a la diversidad de aprendizajes adquiridos durante el desarrollo del presente Trabajo de Fin de Grado, se ha optado por estructurar las conclusiones en tres bloques diferenciados: \textbf{conclusiones científicas}, relacionadas con los principios de rehabilitación y experiencia de usuario; \textbf{conclusiones técnicas}, centradas en el desarrollo del entorno VR; y \textbf{conclusiones personales}, vinculadas al crecimiento profesional y académico del desarrollador.

\subsection*{Conclusiones científicas}

El desarrollo del entorno virtual de rehabilitación ha permitido profundizar en el impacto de la retroalimentación inmediata, la repetición y la interacción activa como elementos clave en el aprendizaje motor. Se ha confirmado que el diseño de ejercicios que exigen movimientos precisos y guiados favorece la consolidación de patrones motores, tal como apuntan estudios previos en neurorehabilitación.

Además, se ha evidenciado la necesidad de adaptar los entornos virtuales a las capacidades físicas y cognitivas del usuario final. El diseño debe contemplar tanto la accesibilidad como la variabilidad del ritmo de rehabilitación de cada paciente. Esto se traduce, por ejemplo, en la necesidad de permitir ajustes de dificultad o controlar el tiempo de sesión.

Finalmente, la experiencia con usuarios reales, aunque limitada, ha servido para validar la utilidad de la realidad virtual como herramienta de apoyo en terapias convencionales, destacando especialmente su potencial motivacional.

\subsection*{Conclusiones técnicas}


Desde el punto de vista técnico, uno de los aprendizajes más significativos ha sido el dominio progresivo del motor \textit{Unreal Engine 5.5}, junto con sus herramientas asociadas como \textit{Blueprints} y el soporte parcial en \textit{C++}. La utilización de la plantilla base para realidad virtual proporcionó una estructura inicial robusta, aunque fue necesario realizar múltiples adaptaciones para responder a los requerimientos específicos del proyecto.

Uno de los principales desafíos fue la creación de un sistema de guardado personalizado capaz de almacenar datos relevantes sobre el rendimiento del usuario. Asimismo, se implementó un mecanismo de exportación de datos en formato \textit{CSV}, utilizando nodos extendidos a través del \textit{Victory Plugin}, lo cual requirió una comprensión profunda del flujo de datos dentro del entorno VR.

También se abordó la configuración dinámica de niveles en función del tipo de ejercicio, lo que implicó una gestión precisa de variables globales, sistemas de navegación y condiciones de activación. La correcta orquestación de los eventos de realidad virtual y la comunicación entre actores mediante \textit{Blueprint Interfaces} y eventos personalizados ha resultado esencial para mantener la coherencia del entorno y garantizar una experiencia fluida.

Por otro lado, se prestó especial atención a la construcción de interfaces accesibles, al tratamiento del input procedente de controladores hápticos, y a la optimización del rendimiento para mantener una tasa de refresco estable en el visor. Todos estos aspectos han contribuido al fortalecimiento del conocimiento sobre la arquitectura de \textit{Unreal Engine}, el ciclo de vida de sus componentes y su modelo de ejecución en tiempo real.


\subsection*{Conclusiones personales}

A nivel personal, este proyecto ha supuesto un reto en todas sus fases, que ha exigido aprender y aplicar conocimientos de programación, diseño de \textit{UX/UI}, metodologías ágiles y documentación técnica. Destacar además de que no se tenia ningún conocimiento sobre el diseño de un videojuego en realidad virtual, esto ha supuesto muchas horas de aprendizaje autodidacta y lectura de documentación.

Se ha aprendido a priorizar tareas, resolver bloqueos técnicos de forma autónoma y a colaborar con un cliente real (ASPAYM), ajustando el desarrollo a sus necesidades y expectativas. Esta experiencia ha resaltado la importancia de la comunicación clara, la empatía con el usuario final y la capacidad de adaptación ante nuevos desafíos.

Por último, este trabajo ha reforzado la idea de que la constancia y la disciplina son más determinantes que la motivación esporádica. La complejidad del desarrollo, junto con la responsabilidad de llevar a cabo un producto funcional en solitario, con todas las fases del desarrollo, ha permitido desarrollar una ética de trabajo sólida y resiliente. Aun así, destacar que la finalidad altruista de este proyecto ha motivado mucho, desde un principio, por sacar el proyecto adelante.

\section{Líneas de trabajo futuras}

El ámbito de la rehabilitación mediante realidad virtual continúa en expansión, y la integración de nuevas dinámicas interactivas, junto con la mejora técnica del sistema, puede marcar una diferencia significativa en su aplicabilidad real. Por ello, se plantean las siguientes líneas de trabajo que permitirían refinar, ampliar y optimizar la solución desarrollada. Estas propuestas se agrupan en diferentes categorías según su naturaleza.

\subsection*{Diseño de nuevos puzles rehabilitadores}

Hasta el momento, se ha implementado únicamente el primer ejercicio interactivo, \textit{El Palacio de Shahriar}, diseñado con prioridad alta por su aplicabilidad en etapas iniciales del proceso de rehabilitación. No obstante, el proyecto contempla la incorporación de un conjunto más amplio de puzles temáticos, que aporten diversidad y progresión en la dificultad. Las siguientes líneas de trabajo están orientadas a completar este objetivo:

\begin{itemize}
    \item \textbf{Desarrollar nuevos puzles:} completar los ejercicios restantes según la lista establecida por ASPAYM.
    \item \textbf{Ajustar la dificultad progresiva:} diseñar cada puzle con niveles de exigencia crecientes, tanto a nivel cognitivo como motriz, adaptados a distintos perfiles de usuario.
    \item \textbf{Incorporar nuevos tipos de interacción:} explorar otras mecánicas más avanzadas (por ejemplo, movimientos finos de muñeca o rotaciones) que complementen el rango de movimientos trabajados.
\end{itemize}

\subsection*{Interfaz y experiencia de usuario}

Aunque la interfaz actual es completamente funcional y ha sido diseñada en base a criterios de accesibilidad, se han detectado oportunidades de mejora a nivel estético y de experiencia de usuario. Algunas líneas a considerar son:

\begin{itemize}
    \item \textbf{Revisión estética:} aplicar un rediseño visual de la interfaz, incorporando una paleta de colores más atractiva, iconografía adaptada y animaciones suaves que refuercen la inmersión.
    \item \textbf{\textit{Feedback} adaptativo:} mejorar el sistema de retroalimentación visual y sonora en función del desempeño del usuario.
    \item \textbf{Personalización de la experiencia:} permitir al usuario configurar preferencias como el volumen o idioma.
\end{itemize}

\subsection*{Rendimiento y optimización}

Pese a los esfuerzos realizados en mejorar el rendimiento durante el desarrollo, es un área que siempre admite margen de mejora, especialmente en entornos de realidad virtual donde los recursos son limitados.

\begin{itemize}
    \item \textbf{Optimización del renderizado:} revisar los modelos y materiales utilizados para garantizar la máxima eficiencia sin comprometer la calidad visual.
    \item \textbf{Reducción del consumo de memoria:} auditar el uso de texturas, sonidos y blueprints para identificar elementos innecesarios o duplicados.
    \item \textbf{Pruebas de rendimiento sistemáticas:} implementar una batería de tests automatizados para comprobar la estabilidad del sistema en diferentes escenarios y niveles.
\end{itemize}

\subsection*{Incorporación de seguimiento de manos (\textit{hand tracking})}

La incorporación de tecnologías de seguimiento de manos puede representar un salto cualitativo en cuanto a accesibilidad, naturalidad del movimiento y eliminación de barreras tecnológicas (como la necesidad de agarrar controladores físicos).

\begin{itemize}
    \item \textbf{Compatibilidad con hand tracking nativo:} adaptar las interacciones del sistema para que puedan ser ejecutadas sin controladores físicos, empleando directamente los gestos naturales de las manos del usuario.
    \item \textbf{Diseño de gestos personalizados:} definir un conjunto de gestos que puedan ser utilizados en sustitución de acciones comunes como agarrar, seleccionar o mover objetos.
    \item \textbf{Evaluación de precisión y robustez:} analizar la fiabilidad del reconocimiento en distintos entornos y condiciones de iluminación, y su impacto en el proceso de rehabilitación.
\end{itemize}


\bibliographystyle{plain}
\bibliography{bibliografia}

\end{document}
